\chapter{PS-PL Bring-Up Guide}

\section*{Purpose}
This chapter documents the first software-driven PS-to-PL bring-up path for the stable wrapper contract on Zynq UltraScale+.

\section{Current Bring-Up Strategy}
The software stack supports two backend modes:
\begin{itemize}[leftmargin=*]
  \item \textbf{fake:} local development and regression without hardware
  \item \textbf{real:} user-space MMIO access intended for Linux on the target board
\end{itemize}

The first real bring-up target is explicitly a `STUB`-mode bitstream.

\section{AXI-Lite Address And Region Expectation}
The real backend assumes software knows the physical base address assigned to the AXI-Lite wrapper in the platform design. The example placeholder used throughout the bring-up documentation is `0xA0000000`, but the actual board value must come from the Vivado address editor and the exported hardware handoff.

The software-visible wrapper region must cover at least `0x1314` bytes, which corresponds to the current register map through the full signature readback window. The backend performs 32-bit little-endian accesses and validates 4-byte alignment for every MMIO transaction.

\section{Real Backend Access Method}
The current real backend uses a `/dev/mem`-style memory-mapped access path suitable for proof-of-concept Linux bring-up. It assumes little-endian 32-bit accesses, page-aligned physical mapping behavior, and a mapped span large enough to cover the documented wrapper register window. The backend handles page-base alignment internally, but the documented wrapper base address from the platform design should still be recorded exactly.

\section{Safe First Sequence}
\begin{enumerate}[leftmargin=*]
  \item Program a bitstream built from `pqsig\_top\_stub\_mode.sv`.
  \item Identify the AXI-Lite wrapper base address from the platform handoff.
  \item Confirm that the assigned wrapper region spans at least `0x1314` bytes.
  \item Configure the daemon for `real` backend mode with that base address.
  \item Run the MMIO probe path to confirm STATUS, ERROR\_CODE, and SIG\_LENGTH visibility.
  \item If needed, clear stale status through the probe path.
  \item Run the explicit STUB selftest only after probe visibility is confirmed.
  \item Confirm the selftest reports `verified\_stub\_signature=true` and `signature\_length=128`.
\end{enumerate}

\section{Engine Mode Selection Note}
The current wrapper uses compile-time or synthesis-time engine-mode selection through the RTL top-level. There is no software-visible runtime mode register. In practice, early board bring-up should use separate bitstreams or explicit build targets for `STUB` and `MLDSA\_OSH` modes.

\section{Transition From STUB to MLDSA\_OSH}
Only after the STUB path is stable on target hardware should the project move to a bitstream configured for `MLDSA\_OSH` mode. The software backend, gRPC API, and MMIO sequencing should remain unchanged across that transition.

\section{Known Bring-Up Blockers}
\begin{itemize}[leftmargin=*]
  \item The current digest-to-message adaptation inside the ML-DSA-OSH shim is provisional.
  \item Key provisioning remains PoC-only static preload data.
  \item Full local mixed-language simulation of the imported core path is still limited by tool availability.
  \item The real backend still depends on actual Zynq address assignment and Linux permissions.
\end{itemize}

\section{Current Claim Boundary}
This phase provides a real executable STUB-mode board procedure. It does not claim that PS-to-PL communication or STUB-mode selftest has already been validated on an actual target board.
